%\documentclass[aspectratio=169]{beamer}
\documentclass{beamer}

\usepackage[utf8]{inputenc}
\usepackage[french]{babel}
\usepackage[page]{appendix}
\usepackage{graphicx}
\graphicspath{{figures/}, {diagrams/}}
\usepackage{tikz}
\usepackage{pgfplots}
\usetikzlibrary{arrows.meta}
\usepackage{smartdiagram}

%\usetheme{tb}
\definecolor{kugreen}{RGB}{66,103,104}
\definecolor{kugreenlys}{RGB}{86,157,160}
\definecolor{kugreenlyslys}{RGB}{165,165,165}
\definecolor{kugreenlyslyslys}{RGB}{242,242,242}
\setbeamercovered{transparent}
\mode<presentation>{%
  \usetheme{PaloAlto}
  \usecolortheme[named=kugreen]{structure}
  \useinnertheme{circles}
  \usefonttheme[onlymath]{serif}
  \setbeamercovered{transparent}
  \setbeamertemplate{blocks}[rounded][shadow=true]
  \addtobeamertemplate{navigation symbols}{}{%
    \usebeamerfont{footline}%
    \usebeamercolor[fg]{footline}%
    \hspace{1em}%
    \insertframenumber/\inserttotalframenumber
  }
}

\title[Service Routière]{\textbf{Implémentation de Service de Collecte et de Stockage des Données sur une plateforme de services Routières}}
%\subtitle{Mémoire de projet de fin d'études}
\author{Moez Bouhlel \and Rihab Majdoub}
\author[Moez B. \and Rihab M.]{\textbf {Moez Bouhlel \and Rihab Majdoub\\[0.2cm] \footnotesize Sous la direction de: \\ Dr. Mohamed Mhiri \and M. Mohamed Amri}}
\institute{Faculté des Sciences de Sfax}
%\titlegraphic{\includegraphics[width=2cm]{logo-djagora}}
\logo{\includegraphics[width=1.6cm]{fss-old}}
\date{05/06/2017}

\begin{document}

\begin{frame}
   \tikz [remember picture,overlay]
    \node at
        ([yshift=-0.8cm, xshift=-1.4cm]current page.north east)
        {\includegraphics[width=2.8cm,height=1.6cm]{logo-djagora}};
   \titlepage
\end{frame}

\begin{frame}
    \frametitle{Plan de la présentation}
    \tableofcontents[hideallsubsections]
\end{frame}

\AtBeginSection[]{%
  \begin{frame}
    \frametitle{Plan de la présentation}
    \tableofcontents[currentsection,hideothersubsections]
  \end{frame}
}

\section{Contexte de projet}

\begin{frame}
    \frametitle{Organisme d'accueil}
    \framesubtitle{Djagora Academy}
\begin{columns}
\begin{column}{0.4\textwidth}
    \begin{figure}
        \includegraphics[width=\textwidth]{logo-djagora.png}
    \end{figure}
\end{column}
\begin{column}{.6\textwidth}
\begin{itemize}
\item Promgramme de mentorat
\item Accélerateur des startups
\end{itemize}
\end{column}
\end{columns}
\end{frame}

\begin{frame}
    \frametitle{Problématique}
    \includegraphics[width=.6\textwidth]{datamanagement-issues}
    \includegraphics[width=.32\textwidth]{question}
\end{frame}

\begin{frame}
    \frametitle{Présentation du projet}
    \framesubtitle{Objectif}
    \begin{columns}
        \begin{column}{0.5\textwidth}
            \begin{figure}
                \includegraphics[width=1\textwidth]{logo-citywatch}
            \end{figure}
        \end{column}
        \begin{column}{0.5\textwidth}
            \begin{itemize}
                \item Startup
                \item 8 membres
                \item Services de consultings
                \item Services routières
            \end{itemize}
        \end{column}
    \end{columns}
\end{frame}

\begin{frame}
    \frametitle{Objectifs du projet}
    \begin{figure}
        \centering
        \includegraphics[width=.73\textwidth]{citywatch-modules}
    \end{figure}
\end{frame}

\section{Gestion projet selon Scrum}

%\begin{frame}
%    \frametitle{Pourquoi Scrum}
%    \centering
%    \includegraphics[width=.7\textwidth]{scrum-logo}
%\end{frame}

\begin{frame}
    \frametitle{Méthode Scrum}
    \centering
    \includegraphics{scrum-model}
\end{frame}

\begin{frame}
\frametitle{Découpage du projet}
\begin{figure}
    \smartdiagram[descriptive diagram]{%
  {Itération 1, {Service de localisation}},
  {Itération 2, {Gestion des rapports}},
  {Itération 3, {Authentification}}}
\end{figure}
\end{frame}

\section{Itération 1: Service Localisation}

\begin{frame}
    \frametitle{But de l'itération}
    \begin{itemize}
        \item<1-> Application mobile qui détecte la position géographique.
        \item<2-> Page web pour afficher la dernière position.
    \end{itemize}
\end{frame}

\begin{frame}
    \frametitle{Diagramme de cas d'utilisation}
    \begin{figure}
        \includegraphics[width=1\textwidth]{sprint1-webservices-usecase}
    \end{figure}
\end{frame}

\begin{frame}
    \frametitle{Diagramme de séquence}
    \begin{figure}
        \includegraphics[width=0.7\textwidth]{sprint1-webservices-post-sequence}
    \end{figure}
\end{frame}

\begin{frame}
    \frametitle{Produit de l'itération}
    \begin{figure}
        \includegraphics[width=0.3\textwidth]{sprint1-android-screenshot}
        \includegraphics[width=0.7\textwidth]{sprint1-dashboard-screenshot}
    \end{figure}
\end{frame}

\begin{frame}
\frametitle{Évaluation de l'itération}
\begin{minipage}{\textwidth}
\usetikzlibrary{plotmarks}

\begin{figure}
\centering
\begin{tikzpicture}[y=.2cm, x=.7cm,font=\sffamily]
\begin{axis}[
xlabel=$Jours$,
ylabel=$Heures\ restantes$,
grid=both,
grid style={line width=.1pt, draw=gray!10},
width=0.8\textwidth,
height=6cm,
%major grid style={line width=.2pt,draw=gray!50},
]
\addplot[color=blue,mark=*] coordinates {
        (0,30)
        (7, 0)
    };
    \addlegendentry{Temps Idéal}

    \addplot[mark=*,red] plot coordinates {
        (0, 30)
        (1, 25)
        (2, 20)
        (3, 12)
        (4, 8)
        (5, 4)
        (6, 0)
        (7, 0)
    };
    \addlegendentry{Temps Restant}
\end{axis}
\end{tikzpicture}
\caption{Graphique d'avancement - Itération 1}
\end{figure}

\end{minipage}
\end{frame}

\section{Itération 2: Gestion des Rapports}

\begin{frame}
    \frametitle{But de l'itération}
    \begin{itemize}
        \item<1-> Déclaration des rapports.
        \item<2-> Détection des secousses et des ralentisseurs.
        \item<3-> Consultation des trajectoires, des rapports, des secousses et des ralentisseurs.
    \end{itemize}
\end{frame}

\begin{frame}
    \frametitle{Diagramme de cas d'utilisation}
    \begin{figure}
        \includegraphics[width=\textwidth]{sprint2-webservices-report-usecase}
    \end{figure}
\end{frame}

\begin{frame}
    \frametitle{Diagramme de séquence}
    \begin{figure}
        \includegraphics[width=.8\textwidth]{sprint2-webservices-report-post-sequence}
    \end{figure}
\end{frame}

\begin{frame}
    \frametitle{Produit de l'itération}
    \begin{figure}
        \includegraphics[width=0.45\textwidth,height=4cm]{sprint2-rapport-screenshot1}
        \includegraphics[width=0.5\textwidth]{sprint2-dashboard-screenshot2}
    \end{figure}
\end{frame}

\section{Itération 3: Authentification}

\begin{frame}
    \frametitle{But de l'itération}
    \begin{itemize}
        \item<1-> Détecte et projection de vitesse moyenne, d'état du réseau cellulaire.
        \item<2-> Authentification.
        \item<3-> Business Intelligence.
    \end{itemize}
\end{frame}

\begin{frame}
    \frametitle{Diagramme de cas d'utilisation}
    \begin{figure}
        \includegraphics[width=\textwidth]{sprint3-webservices-oauth-usecase}
    \end{figure}
\end{frame}

\begin{frame}
    \frametitle{Diagramme de séquence}
    \begin{figure}
        \includegraphics[width=.65\textwidth]{sprint3-webservices-oauth-sequence}
    \end{figure}
\end{frame}

\begin{frame}
    \frametitle{Produit de l'itération}
    \begin{figure}
        \includegraphics[width=0.45\textwidth,height=4cm]{sprint2-rapport-screenshot1}
        \includegraphics[width=0.5\textwidth]{sprint2-dashboard-screenshot2}
    \end{figure}
\end{frame}


\section{Réalisation}

\begin{frame}
    \frametitle{Technologies utilisées}
    \begin{itemize}
        \item Architecture MVC
        \item Architecture RESTful
        \item Format JSON
        \item PHP (Laravel) + MySQL
        \item JavaScript + HTML5
        \item Java + Android
    \end{itemize}
\end{frame}

\begin{frame}
    \begin{center}
    \begin{figure}
        \includegraphics[width=0.9\textwidth]{demo}
    \end{figure}
    \end{center}
\end{frame}

\section{Conclusion \& Perspectives}

\begin{frame}
    \frametitle{Conclusion}
    \begin{itemize}
        \item Développer des Services Web tel que: service de position et service de ralentisseurs.
        \item Découper le projet suivant la méthode Scrum en 3 itérations: service de localisation, gestion des rapport et authentification.
        \item Utiliser des technologies modernes comme JSON, RESTful, Android.
    \end{itemize}
\end{frame}

\begin{frame}
    \frametitle{Perspectives}
    \begin{itemize}
        \item Appliquer la solution dans des scénarios comme calcule du chemin optimal suivant les caractéristiques des marchandises.
        \item Étendre l'application dans un environnement du Big Data comme Business Intelligence, data mining.
        \item Valoriser notre solution pour l'ouverture d'une startup de service de consulting.
    \end{itemize}
\end{frame}

\begin{frame}
    \begin{center}
        \bfseries \Huge
        Merci
    \end{center}
\end{frame}

\section*{Annexe}
\begin{frame}
\frametitle{Diagramme 1}
\end{frame}

\end{document}
